(1) Schrödinger方程外域满足 \(\psi''=\kappa^2\psi\)，且 \(\psi(\pm a^+)=\psi(\pm a^-)\)。跨越每个吸引Delta点积分给 \(\psi'(x_0^+)-\psi'(x_0^-)=-2mg\psi(x_0)/\hbar^2\)。定义 \(\alpha=mg/\hbar^2>0\)，两阱耦合必须同时满足两个跳跃条件。

(2) 可用偶解 \(\psi_+=C_+[e^{-\kappa|x-a|}+e^{-\kappa|x+a|}]\) 与奇解 \(\psi_-=C_-[e^{-\kappa|x-a|}-e^{-\kappa|x+a|}]\)。代入 \(x=\pm a\) 的导数跳跃，得到 \(\boxed{\kappa_+=\alpha(1+e^{-2\kappa_+a})}\)、\(\boxed{\kappa_-=\alpha(1-e^{-2\kappa_-a})}\)，前者总有唯一正根。

(3) 奇态方程在 \(\kappa=0\) 有零形式根，但它不是可归一的束缚态。比较右侧在原点的斜率 \(2\alpha a\) 与左侧斜率1，只有 \(\boxed{2\alpha a>1}\) 才出现严格正的奇态根。\(a\to\infty\) 时两个态均趋 \(\kappa=\alpha\)；\(a\to0\) 时偶态趋 \(\kappa=2\alpha\)，对应合并后强度 \(2g\) 的单Delta阱，奇态消失。

(4) 定义 \(f_\pm(x)=e^{-\kappa_\pm|x-a|}\pm e^{-\kappa_\pm|x+a|}\)。自项积分总共 \(2/\kappa\)，交叉重叠积分为 \(e^{-2\kappa a}(2a+1/\kappa)\)，故 \(\boxed{C_\pm^{-2}=2/\kappa_\pm\pm2e^{-2\kappa_\pm a}(2a+1/\kappa_\pm)}\)。势阱中心反射对称，奇偶波函数乘积为奇函数，因此两态正交。
